\chapter{Kennfelder aus Feldberechnungen}
\label{kap:format-kennfelder}

Ein Magnet, ein Motor oder eine Drossel mit Eisen lässt sich nicht mit einer
Formel beschreiben. Ein Feldprogramm (FEM) rechnet ihn, \esim\ rechnet die
Schaltung und die Mechanik drumherum. Der einfachste Weg, beide zu koppeln, ist
ein \emph{Kennfeld}: Das Feldprogramm rechnet vorab für ein Gitter aus Strom
und Position die Flussverkettung und die Kraft, \esim\ liest die Werte als
Tabelle über zwei Größen (Abschnitt~\ref{sec:tabellen}) und rechnet damit
Schaltung und Bewegung. Das Feldprogramm läuft dabei nicht mit; es braucht
keine Schnittstelle zu \esim.

Das Kapitel beschreibt den Ablauf an einem Topfmagneten mit Flachanker:
Beispiele \datei{examples/hubmagnet_blocked.json} und
\datei{examples/hubmagnet_motion.json}, Daten
\datei{examples/data/hubmagnet_map.csv}. Das Kennfeld stammt aus dem
FEM-Programm ITEFEM (3D, Kantenelemente, nichtlineares Eisen, 400 Windungen,
Hub 2 bis 8\,mm).

\section{Wann ein Kennfeld, wann ein Plugin}
\label{sec:kennfeld-oder-plugin}

Ein Kennfeld ersetzt das Feldprogramm durch seine statische Antwort. Es passt,
wenn

\begin{itemize}
  \item das Feld der Bewegung und dem Strom ohne Verzögerung folgt
    (quasistatisch, keine Wirbelströme, keine Hysterese),
  \item das Feldprogramm nicht als Bibliothek aufrufbar ist oder je Punkt
    Minuten braucht,
  \item der Zustand des Bauteils durch wenige Größen (Strom, Position)
    bestimmt ist.
\end{itemize}

Wirbelströme, Hysterese und Zustände mit eigener Zeitkonstante verlangen ein
Plugin (Kapitel~\ref{kap:format-plugins}); dort rechnet die Bibliothek bei
jeder Auswertung mit. Die Aufteilung in Schaltung, Mechanik und Feld bleibt
dieselbe, nur die Quelle von $\psi$ und $F$ wechselt.

\section{Ablauf}
\label{sec:kennfeld-ablauf}

\begin{enumerate}
  \item \textbf{Größen und Gitter festlegen} (Abschnitte~\ref{sec:kennfeld-groessen}
    und~\ref{sec:kennfeld-gitter}).
  \item \textbf{Rechnen:} das Feldprogramm für jeden Gitterpunkt, eine Datei mit
    einem Punkt je Zeile (Abschnitt~\ref{sec:kennfeld-datei}).
  \item \textbf{Prüfen} mit \datei{tools/check_map.py}
    (Abschnitt~\ref{sec:kennfeld-pruefen}).
  \item \textbf{Einbinden:} Tabelle, Flusselement, Mechanik
    (Abschnitt~\ref{sec:kennfeld-einbinden}).
  \item \textbf{Prüfen gegen eine Rechnung des Feldprogramms} mit gleicher
    Ansteuerung (Abschnitt~\ref{sec:kennfeld-pruefen-lauf}).
\end{enumerate}

\section{Größen und Vorzeichen}
\label{sec:kennfeld-groessen}

\begin{longtable}{@{}P{2.6cm}P{11.2cm}@{}}
  \toprule
  Größe & Bedeutung \\
  \midrule
  \endhead
  \bottomrule
  \endfoot
  \code{i} & Strom der Wicklung in A. Bei mehreren Wicklungen in Reihe der gemeinsame Strom \\
  \code{g} & Position in m. Im Beispiel der Luftspalt zwischen Kern und Anker, nicht die Lage des Ankers im Raum \\
  \code{psi} & Flussverkettung der \emph{ganzen} Spule in Vs, also mit der Windungszahl \\
  \code{F} & Betrag der Anziehungskraft in N, positiv \\
\end{longtable}

\begin{itemize}
  \item Alles in SI-Einheiten. Das Feldprogramm rechnet oft in mm; die Umrechnung
    gehört in die Datei, nicht in die Netzliste.
  \item Das Vorzeichen der Kraft legt die Mechanik fest, nicht das Kennfeld:
    $F$ ist ein Betrag, die Richtung steht in der Bewegungsgleichung
    (Abschnitt~\ref{sec:kennfeld-einbinden}).
  \item $\psi$ ist ungerade und $F$ gerade im Strom. Man rechnet daher nur
    $i\ge0$ und spiegelt: $\psi(-i,g)=-\psi(i,g)$, $F(-i,g)=F(i,g)$. Die Datei
    muss trotzdem das volle Gitter mit beiden Vorzeichen enthalten.
  \item Bei $i=0$ ist $\psi=0$ und $F=0$ (kein Dauermagnet, keine Remanenz). Der
    Punkt braucht nicht gerechnet zu werden.
\end{itemize}

\section{Dateiformat}
\label{sec:kennfeld-datei}

Eine Textdatei mit Kopfzeile, ein Punkt je Zeile, Komma als Trenner und Punkt
als Dezimalzeichen:

\begin{shell}
i,g,psi,F,F_W
-8,0.002,-0.138774124,111.477,112.907482
-8,0.003,...
...
 8,0.008,...
\end{shell}

\begin{itemize}
  \item Die Reihenfolge der Zeilen ist beliebig. Jede Kombination der Stützwerte
    von \code{i} und \code{g} kommt \emph{genau einmal} vor (Gitter 17 mal 6
    im Beispiel: 102 Zeilen).
  \item Die Achsen müssen nicht gleichmäßig geteilt sein, aber streng steigen.
  \item Zeilen, die nicht mit einer Zahl beginnen, überspringt \esim; die
    Kopfzeile darf also stehen.
  \item Mehrere Größen in einer Datei sind üblich: \esim\ legt für jede Spalte
    eine eigene Tabelle an (\code{zcol}), siehe unten. Spalten, die nicht
    gebraucht werden, schaden nicht. Im Beispiel steht neben $F$ (Maxwell-Tensor)
    die Kraft \code{F_W}, die aus der Koenergie von $\psi$ folgt.
\end{itemize}

\section{Gitter wählen}
\label{sec:kennfeld-gitter}

Die Interpolation ist kubisch (Hermite, stetige Steigung); zwischen den
Stützstellen sieht das Newton-Verfahren glatte Kurven, aber nur so genau, wie
die Stützstellen den Verlauf tragen. Daraus folgt:

\begin{itemize}
  \item \textbf{Betriebsbereich abdecken.} Außerhalb des Gitters hält \esim\ den
    Randwert (\code{"extrap": "hold"}). Fährt die Bewegung über den Rand hinaus,
    rechnet es dort mit falscher Kraft, ohne Meldung. Anschläge in der Mechanik
    (Abschnitt~\ref{sec:kennfeld-einbinden}) halten die Position im Gitter.
  \item \textbf{Dichter, wo es knickt.} Bei Sättigung ändert sich $\psi(i)$ über
    wenige Stützstellen stark. Im Beispiel ist $\psi(8\,\mathrm{A})$ bei 2\,mm
    nur 18\,\% größer als $\psi(6\,\mathrm{A})$; die Stützstellen 6 und 8\,A
    erfassen den Knick nur grob. Ströme mit geeigneter Staffelung (1, 2, 3, 4, 6, 8\,A)
    sind besser als gleiche Abstände.
  \item \textbf{Kleine Ströme.} Die Induktivität steigt bei sehr kleinen Strömen
    (Anfangspermeabilität des Eisens). Fehlen Stützstellen zwischen 0 und dem
    ersten Strom, liegt die Anfangssteigung falsch: im Beispiel 15,3\,mH statt
    14,4\,mH, mit Stützstellen bei 0,25 und 0,5\,A behoben. Das wirkt auf schnelle
    Schaltvorgänge am Anfang.
  \item \textbf{Aufwand.} Jeder Punkt ist eine Feldrechnung (im Beispiel 4 bis
    10 Minuten). Man rechnet die Stützstellen in der Reihenfolge ihrer
    Wichtigkeit und prüft das Gitter, bevor man es verfeinert.
\end{itemize}

\section{Prüfen}
\label{sec:kennfeld-pruefen}

\datei{tools/check_map.py} prüft eine Kennfelddatei, bevor sie in eine
Rechnung geht (Python mit numpy und scipy):

\begin{shell}
python3 tools/check_map.py examples/data/hubmagnet_map.csv
python3 tools/check_map.py map.csv --x i --y g --psi psi --force F --tol 5
\end{shell}

\begin{longtable}{@{}P{2.6cm}P{11.2cm}@{}}
  \toprule
  Prüfung & Inhalt \\
  \midrule
  \endhead
  \bottomrule
  \endfoot
  Gitter & jede Kombination genau einmal, kein doppelter Punkt, beide Achsen steigend \\
  Null & $\psi=0$ und $F=0$ bei $i=0$ \\
  Symmetrie & $\psi$ ungerade, $F$ gerade, wenn die Datei negative Ströme enthält \\
  Monotonie & $\psi$ steigt mit dem Strom \\
  Energie & die Kraft folgt aus der Flussverkettung (siehe unten) \\
\end{longtable}

\paragraph{Energieverträglichkeit.}
Strom und Position sind die Zustandsgrößen; die Koenergie
$W'(i,g)=\int_0^i\psi\,\mathrm{d}i'$ ergibt die Kraft
$F=-\partial W'/\partial g$. Wird $F$ unabhängig berechnet (Maxwell-Tensor),
stimmt sie nie genau mit dieser Ableitung überein. Das Werkzeug vergleicht beide
je Position. Für das Beispiel:

\begin{longtable}{@{}rl@{}}
  \toprule
  Luftspalt & $F_W$ gegen $F$ (für $i\ge0{,}8$\,A) \\
  \midrule
  \endhead
  \bottomrule
  \endfoot
  2\,mm & $+1{,}2$ bis $+1{,}8$\,\% \\
  3\,mm & $+0{,}2$ bis $+1{,}2$\,\% \\
  4\,mm & $+4{,}3$ bis $+4{,}7$\,\% \\
  5\,mm & $+6{,}6$ bis $+6{,}9$\,\% \\
  6\,mm & $+6{,}3$ bis $+7{,}9$\,\% \\
  8\,mm & $+9{,}4$ bis $+13{,}8$\,\% \\
\end{longtable}

Die Abweichung wächst mit dem Luftspalt; im Beispiel ist die Schicht des
Maxwell-Tensors bei großem Spalt dünner als ein Element des Netzes. Die Wahl:

\begin{itemize}
  \item Für die Rechnung die Kraft aus der Koenergie nehmen (\code{F_W}). Sie
    ist zu $\psi$ energiekonsistent: Die elektrische Energie, die über
    $\mathrm{d}\psi/\mathrm{d}t$ in den Magneten fließt, und die mechanische
    Arbeit gehen auf.
  \item Die unabhängig berechnete Kraft $F$ als Kontrolle behalten. Im Beispiel
    verschiebt sie den Zeitpunkt des Anzugs um 3\,\% (17,3 gegen 17,9\,ms); der
    Maximalstrom ist gleich.
  \item Weicht die Kraft um mehr als wenige Prozent ab, die Stützstellen mit
    Abweichung nachrechnen (feinere Schicht oder Netz) oder die Ursache im
    Feldprogramm suchen.
\end{itemize}

Bei Fehlern im Gitter meldet \esim\ schon beim Einlesen, etwa

\begin{shell}
error: table "Fmag": the grid has no point at (-8.000000, 0.002000); 17 x 6 values are required
error: table "Fmag": the point (-8.000000, 0.005000) appears twice
error: table "Fmag": cannot open 'nofile.csv'
\end{shell}

\section{Einbinden}
\label{sec:kennfeld-einbinden}

Das Beispiel \datei{examples/hubmagnet_blocked.json} hält den Anker bei
$g=4$\,mm fest und legt 12\,V an:

\begin{json}
"parameters": { "R": 1.5, "U": 12, "g": "4m" },
"tables": {
  "psi":  { "file": "data/hubmagnet_map.csv", "xcol": 0, "ycol": 1, "zcol": 2, "interp": "cubic" },
  "Fmag": { "file": "data/hubmagnet_map.csv", "xcol": 0, "ycol": 1, "zcol": 4, "interp": "cubic" }
},
"components": [
  { "type": "V", "name": "V1", "nodes": ["a", "0"],
    "wave": { "type": "pulse", "v1": 0, "v2": "U", "delay": 0, "rise": "1n" } },
  { "type": "R", "name": "R1", "nodes": ["a", "m"], "value": "R" },
  { "type": "B", "name": "B1", "nodes": ["m", "0"], "phi": "psi(i(B1), g)" }
],
"variables": { "F": { "guess": 0 } },
"equations": [ "F = Fmag(i(B1), g)" ]
\end{json}

\begin{itemize}
  \item Jede Spalte der Datei wird eine eigene Tabelle: \code{xcol} und \code{ycol}
    sind Strom und Position, \code{zcol} die Spalte des Wertes (gezählt ab 0).
  \item Die Spule ist ein Flusselement \code{B} mit
    \code{"phi": "psi(i(B1), g)"} (Abschnitt~\ref{sec:bquelle}); die Spannung ist
    $\mathrm{d}\psi/\mathrm{d}t$. Der Widerstand der Wicklung ist ein eigenes
    Bauteil \code{R}, er steht nicht im Kennfeld.
  \item \esim\ bildet die Ableitungen $\partial\psi/\partial i$ und
    $\partial\psi/\partial g$ aus der Tabelle (exakt zur Interpolation). Im
    Strom gibt es die inkrementelle Induktivität, in der Position die
    Bewegungsspannung: Ist \code{g} eine Variable der Bewegungsgleichung, enthält
    $\mathrm{d}\psi/\mathrm{d}t$ den Term $\partial\psi/\partial g\cdot\dot g$.
    Die Wechselstromanalyse linearisiert um den Arbeitspunkt.
  \item Die Kraft liefert die zweite Tabelle. Sie geht als Variable in die
    Mechanik.
\end{itemize}

\subsection*{Freier Anker}

\datei{examples/hubmagnet_motion.json} lässt den Anker beweglich. Die Position ist
eine Variable, die Mechanik besteht aus Masse, Feder, Dämpfer und Anschlägen:

\begin{json}
"parameters": { "m": 0.05, "k": 500, "c": 5, "g0": "8m",
                "gmin": "2m", "gmax": "8m", "kc": 1e6, "cc": 200 },
"variables": { "g": { "init": "gmax" }, "v": { "init": 0 } },
"equations": [
  "der(g) = v",
  "m*der(v) = -Fw(i(B1), g) + k*(g0 - g) - c*v + kc*max(0, gmin - g) - kc*max(0, g - gmax) - cc*v*(g < gmin)"
]
\end{json}

\begin{itemize}
  \item \code{g} ist die Position, die das Kennfeld als Eingang bekommt. Die Kraft
    des Magneten zieht den Luftspalt zu: $-F_W$ in der Gleichung.
  \item Die Anschläge sind steife Federn (\code{kc}) mit Dämpfung (\code{cc}),
    die nur im Anschlag wirken. Sie halten den Anker im Bereich des Kennfeldes,
    hier 2 bis 8\,mm. Ohne sie würde \esim\ außerhalb den Randwert halten.
  \item Der Verlauf: Nach etwa 17\,ms ist der Anker angezogen. Der Strom bricht
    während der Bewegung ein (Bewegungsspannung), steigt danach auf 8\,A; nach dem
    Abschalten bei 61\,ms öffnet die Feder den Magneten wieder.
\end{itemize}

\begin{hinweis}
  Das Kennfeld ersetzt keine Ereignisse: Schlägt der Anker an, ist das ein
  steifer Kontakt in den Gleichungen, kein Sprung in der Geschwindigkeit. Mit
  einer Anschlagsteifigkeit $10^6$\,N/m und der kleinen Dämpfung bleibt die
  Schrittweite des Integrators dabei im Rahmen; sehr steife Anschläge machen
  die Rechnung langsam.
\end{hinweis}

\section{Prüfen gegen das Feldprogramm}
\label{sec:kennfeld-pruefen-lauf}

Ein Kennfeld ist erst geprüft, wenn eine Rechnung damit ein Ergebnis liefert, das
das Feldprogramm selbst gegen die Zeit erzeugt. Der Ablauf für das Beispiel:

\begin{enumerate}
  \item Das Feldprogramm rechnet den Fall \emph{gekoppelt} mit der Schaltung:
    Spule mit Widerstand 1,5\,\ohm\ an einer Spannungsquelle (Sprung auf
    12\,V), Anker festgehalten bei 4\,mm, Zeitschritte aus einer vorgegebenen
    Liste. Ergebnis: $t$, $i$, $\psi$.
  \item \esim\ rechnet denselben Fall aus dem Kennfeld
    (\datei{examples/hubmagnet_blocked.json}).
  \item Man vergleicht $i(t)$ und $\psi(t)$.
\end{enumerate}

Beim Vergleich muss man zwei Fehlerquellen trennen: das Kennfeld (Gitter,
Interpolation) und die Zeitdiskretisierung der beiden Programme. Dafür löst man
die Gleichung $\mathrm{d}\psi/\mathrm{d}t = U-Ri$ mit $\psi(i)$ aus dem
Kennfeld \emph{genau} (zum Beispiel mit einem Python-Löser) und vergleicht beide
Rechnungen damit. Im Beispiel:

\begin{longtable}{@{}P{7.6cm}P{6.2cm}@{}}
  \toprule
  Rechnung & größte Abweichung im Strom \\
  \midrule
  \endhead
  \bottomrule
  \endfoot
  \esim\ (BDF mit Fehlerkontrolle, Kennfeld als Tabelle) & 0,09\,mA \\
  Feldprogramm (Rückwärts-Euler, 60 Schritte) & 160\,mA (2\,\%) \\
\end{longtable}

\esim\ liegt auf der Lösung des Kennfeldes; die Abweichung des Feldprogramms ist
die Zeitdiskretisierung bei 60 Schritten (am größten bei 12 bis 25\,ms, wo der
Strom steil steigt), nicht das Kennfeld. Fehlt ein feineres Zeitraster, rechnet
man \esim\ auch mit dem Rückwärts-Euler-Verfahren auf derselben Zeitliste
(\code{"method": "be"}, Kapitel~\ref{kap:format-analysen}) und vergleicht
das; dann bleibt nur die Wirkung des Kennfeldes übrig.

\section{Grenzen}
\label{sec:kennfeld-grenzen}

\begin{itemize}
  \item \textbf{Statisch.} Das Kennfeld kennt keine Wirbelströme, keine Hysterese,
    keine Verzögerung des Feldes. Für Wechselströme hoher Frequenz oder massives
    Eisen taugt es nicht.
  \item \textbf{Zwei Achsen.} Strom und Position. Mehrere Wicklungen, Winkel und
    Position zugleich oder zusätzlich die Temperatur brauchen mehrere
    Tabellen oder ein Plugin.
  \item \textbf{Rand.} Außerhalb des Gitters hält \esim\ den Randwert. Der
    Betriebsbereich muss im Gitter liegen (Anschläge).
  \item \textbf{Nicht monoton.} Die kubische Interpolation über zwei Größen kann
    zwischen groben Stützstellen überschwingen. Ein feineres Gitter und die
    Prüfung der Monotonie (\datei{tools/check_map.py}) fangen das ab.
  \item \textbf{Glatte Daten.} Rauschen in den Feldrechnungen (Netzwechsel, Kantenwechsel
    der Kraft) stört das Newton-Verfahren. Eine Kennfelddatei sollte aus
    glatten, konvergierten Rechnungen stammen.
\end{itemize}
